% Source: https://artofproblemsolving.com/community/c3h2523504_a_problem
% Stable identity: aops-topic-2523504-q01
Greg's number has expected value $\dfrac{a+1}{2}$ and Nikhil's $\dfrac{b+1}{2}$. By independence the expected product is the product of the expected values, so
\[\dfrac{(a+1)(b+1)}{4}=3\left(\dfrac{a+1}{2}+\dfrac{b+1}{2}\right).\]
With $A=a+1\ge2$ and $B=b+1\ge2$ this is $AB=6A+6B$, i.e.
\[(A-6)(B-6)=36.\]
If $A-6<0$, then $A-6\in\{-4,\dots,-1\}$, so $B-6=\dfrac{36}{A-6}\le-9$ and $B<0$, impossible. So $A-6$ is a positive divisor of $36$:
\[A-6\in\{1,2,3,4,6,9,12,18,36\},\qquad a=A-1\in\{6,7,8,9,11,14,17,23,41\},\]
each with a positive integer $b$. The sum is
\[6+7+8+9+11+14+17+23+41=\boxed{136}.\]
